\documentclass[10pt]{article}
\usepackage{tcolorbox}
\tcbuselibrary{theorems}

\def\Type{LessonCaelan}
\def\TypeNum{Lesson 19 Key}
\def\TypeTitle{Summations}
\def\Semester{Fall 2021} %Not Used 
\def\Course{MATH 1190 \& 1210}
\def\Targets{  \bf\color{Fuchsia}I3, \bf\color{Fuchsia}I4 }

\usepackage{preambleDJ}

%%%%%%%%%%%% BEGIN DOCUMENT %%%%%%%%%%%%%%%
%%%%%%%%%%%% BEGIN DOCUMENT %%%%%%%%%%%%%%%
%%%%%%%%%%%% BEGIN DOCUMENT %%%%%%%%%%%%%%%
%%%%%%%%%%%% BEGIN DOCUMENT %%%%%%%%%%%%%%%
%%%%%%%%%%%% BEGIN DOCUMENT %%%%%%%%%%%%%%%

%%%%%%%%%%%% BEGIN DOCUMENT %%%%%%%%%%%%%%%
%%%%%%%%%%%% BEGIN DOCUMENT %%%%%%%%%%%%%%%
%%%%%%%%%%%% BEGIN DOCUMENT %%%%%%%%%%%%%%%
%%%%%%%%%%%% BEGIN DOCUMENT %%%%%%%%%%%%%%%
%%%%%%%%%%%% BEGIN DOCUMENT %%%%%%%%%%%%%%%

\begin{document}
\pagestyle{otherpages}
\thispagestyle{firstpage}
\vs{.3}

\textbf{Intended Learning Outcomes}
After this lesson, you will be able to 

\begin{itemize}
    \item Compute expressions involving summation notation. (\target{F3})
    \item Write sums concisely using summation notation. (\target{F3})
    \item Understand and apply the linearity properties of summation (\target{F3})
    %\item Use summation properties to expand and simplify sums (\target{F3})
\end{itemize}

\vspace{-0.6cm}

\hrulefill\\

%Some important properties of probability rely on being able to add several, similarly-formatted quantities together. For example, one important property of probability is that when we sum together the probabilities of each outcome in the sample space, we should get 1 since there is a 100\% chance that {\it something} in the sample space will happen:
%\[ \text{If }S=\{outcome_1,outcome_2,\ldots,outcome_n\},\text{ then }P(\{outcome_1\}) + P(\{outcome_2\}) + \ldots + P(\{outcome_n\}) = 1.\]
%As we'll see in the coming lessons, the above statement can be more succinctly written as a summation:
%\[ \sum_{i=1}^{n} P(\{outcome_n\}) = 1 \quad\text{ or }\quad \sum_{outcome\in S} P(outcome) = 1\]
In this class, we'll practice using notation for expressing sums using summation so that we will be able to use it in our upcoming probability lessons.
\important{Anatomy of a Summation}
{
\[
\sum_{i=1}^{n} a_i
\]

\[
\begin{array}{lcl}
  &\text{Summation symbol (means “add up”)}\\[6pt]
i &\text{Index of summation (variable that changes)}\\[6pt]
1 & \text{Starting value (first $i$ used)}\\[6pt]
n & \text{Ending value (last $i$ used)}\\[6pt]
a_i &  \text{Summand (expression being added)}
\end{array}
\]
}
\vfill

{\example / \exercise}\ltrm{\bf\color{ProcessBlue}F3} Expand and compute the following sums:
    \begin{enumerate}
    \item[(a)] ({\bf\color{blue}Example}) $\displaystyle\sum_{i=0}^2 \frac{2i+1}{i+1}$
        \par\noindent\textcolor{blue}{\emph{Solution:} Evaluate termwise: for $i=0,1,2$ we get $1,\;3/2,\;5/3$. Thus
        \[
        \sum_{i=0}^2 \frac{2i+1}{i+1}=1+\frac32+\frac53=\frac{25}{6}.
        \]}
        \vfill
    \item[(b)] ({\bf\color{blue}Exercise}) $\displaystyle\sum_{j=1}^4 \log_2\left( 2j\right)$
        \par\noindent\textcolor{blue}{\emph{Solution:} For $j = 1,2,3,4$
        \begin{align*}
        \log_2(2\cdot 1)&= \log_2(2) = 1, \text{ ($j = 1$) }\\
        \log_2(2\cdot 2) &= \log_2(4) = 2, \text{ ($j = 2$) }\\
        \log_2(2\cdot 3) &= \log_2(2) + \log_2(3)\\
        &= 1 + \log_2(3) \text{ ($j = 3$),  and }\\
        \log_2(2\cdot 4) &= \log_2(2) + \log_2(4)= 1 + 2 \text{ ($j = 4$) }
        \end{align*}
        Putting these all together, we've 
        $$\displaystyle\sum_{j=1}^4 \log_2\left( 2j\right) = 1 + 2 + (1 + \log_2(3)) + (1 + 2) = 7 + \log_2(3)$$
        Alternatively, we've 
        \begin{align*}
        \sum_{j=1}^4 \log_2(2j)&= \sum_{j=1}^4 (\log_2 2 +\log_2 j)\\
        &= \sum_{j=1}^4 (1+\log_2 j)\\
        &= (1+\log_2 1) + (1 + \log_2 2) + (1 + \log_2 3) + (1 + \log_2 4)\\
        &= (1 + 0) + (1 + 1) + (1 + \log_2 3) + (1 + 2)\\
        &= 7+\log_2 3\\
        \end{align*}
        So the sum is $7+\log_2 3$. }
        \vfill
        
    \item[(c)] ({\bf\color{blue}Exercise}) $\displaystyle\sum_{k=0}^3 k(-1)^k$
        \par\noindent\textcolor{blue}{\emph{Solution:} Terms: 
        \begin{align*}
            0\cdot (-1)^{0} &= 0 \text{ $(k = 0)$ }\\
            1\cdot (-1)^{1} &= -1 \text{ $(k = 1)$ }\\
            2\cdot (-1)^{2} &= 2 \text{ $(k = 2)$ }\\
            3\cdot (-1)^{3} &= -3 \text{ $(k = 3)$ }\\
        \end{align*}
        Sum $0-1+2-3=-2$. So, $\displaystyle\sum_{k=0}^3 k(-1)^k = -2$.}
        \vfill
    \end{enumerate}
    
{\example / \exercise}\ltrm{\bf\color{ProcessBlue}F3} Write - but don't compute - the following sums using summation notation:
    \begin{enumerate}
    \item[(a)] ({\bf\color{blue}Example}) $2+3+4+5+6$
        \par\noindent\textcolor{blue}{\emph{Solution:} $\displaystyle\sum_{n=2}^{6} n$.
        Alternatively, $\displaystyle\sum_{n=1}^{5} n+ 1$
        
        
        }
        \vfill
    \item[(b)] ({\bf\color{blue}Exercise}) $\displaystyle\frac12+\frac14+\frac16+\frac18+\frac1{10}+\frac1{12}$
        \par\noindent\textcolor{blue}{\emph{Solution:} These are the reciprocals of the even integers $2,4,\dots,12$, so one compact form is
        \[
        \sum_{k=1}^{6}\frac{1}{2k} \quad\text{(equivalently }\frac{1}{2}\sum_{k=1}^{6}\frac{1}{k}\text{).}
        \]}
        \vfill
    \item[(c)] ({\bf\color{blue}Exercise}) $\displaystyle-\frac15+\frac2{25}-\frac3{125}+\frac4{625}$
        \par\noindent\textcolor{blue}{\emph{Solution:} For $k=1,\dots,4$ the $k$-th term is $(-1)^k\frac{k}{5^k}$, so
        \[
        -\frac15+\frac{2}{25}-\frac{3}{125}+\frac{4}{625}=\sum_{k=1}^4 (-1)^k\frac{k}{5^k}.
        \]
        Think about these in terms of the PATTERNS you see between the terms. For example, in (b):
        \begin{itemize}
            \item The terms are always positive\\
            \item There is always a 1 in the numerator\\
            \item The denominator is always a multiple of 2.
        \end{itemize}
        In c):
        \begin{itemize}
            \item The signs of the terms are alternating (-  to + )\\
            \item The natural numbers 1 to 4 occur in the numerator.\\
            \item Powers of 5 appear in the denominator.
        \end{itemize}
        }
        \vfill
    \end{enumerate}
    
\newpage

\important{Summation Properties}
{
Let $c$ be any constant, and let $m$ and $n$ be whole numbers with $m<n$.
\begin{itemize}
\item $\displaystyle\sum_{i=m}^n(a_i+b_i) = \sum_{i=m}^n a_i + \sum_{i=m}^n b_i$ \hfill (linearity property \#1)
\item $\displaystyle\sum_{i=m}^n ca_i = c\sum_{i=m}^n a_i$  \hfill (linearity property \#2)
\item $\displaystyle\sum_{i=1}^n c = cn$
\end{itemize}
}
\par\noindent\textcolor{blue}{
\begin{itemize}
    \item 
\end{itemize}

{\example}\ltrm{\bf\color{ProcessBlue}F3} Use the properties of summations exhaustively to decompose the following summation: \[\sum_{j=1}^n (a_j-3b_j+7c_j-4)\]

    \par\noindent\textcolor{blue}{\emph{Solution:} Apply linearity termwise:
    \begin{align*}
    \sum_{j=1}^n (a_j-3b_j+7c_j-4) &=\sum_{j=1}^n a_j + \sum_{j=1}^n -3b_j +\sum_{j=1}^n 7c_j + \sum_{j=1}^n - 4 \text{ (by linearity property 1) }\\
    &=\sum_{j=1}^n a_j -3\sum_{j=1}^n b_j +7\sum_{j=1}^n c_j - \sum_{j=1}^n 4 \text{ (by linearity property 2) }\\
    \text{ and } \sum_{j=1}^n 4 &=4n \text{ (by linearity property 3). }
    \end{align*}
    So,
    \[
    \sum_{j=1}^n (a_j-3b_j+7c_j-4) =\sum_{j=1}^n a_j -3\sum_{j=1}^n b_j +7\sum_{j=1}^n c_j -4n.
    \]}
    
    \vfill
    
{\exercise}\ltrm{\bf\color{ProcessBlue}F3} Use the properties of summations exhaustively to decompose the following summations:
    \begin{enumerate}
    \item[(a)] $\displaystyle\sum_{i=1}^n (2a_i+3b_i-1)$
        \par\noindent\textcolor{blue}{\emph{Solution:} Linearity gives
        \[
        \sum_{i=1}^n (2a_i+3b_i-1)=2\sum_{i=1}^n a_i +3\sum_{i=1}^n b_i -\sum_{i=1}^n 1
        =2\sum_{i=1}^n a_i +3\sum_{i=1}^n b_i - n.
        \]}
        \vfill
    \item[(b)] $\displaystyle\sum_{k=1}^n \left(\frac{a_k\cdot k}4+9b_k^2-2c_k\cdot d_k\right)$\quad(\emph{Hint:} be careful - this is a bit of a trap.)
        \par\noindent\textcolor{blue}{\emph{Solution:} Be careful: $k$ is not a constant with respect to the summation index, so it cannot be pulled outside. \textbf{Note:} Sums do NOT distribute across multiplication or division. Anything depending on the index cannot be factored out! Applying linearity termwise yields
        \[
        \sum_{k=1}^n \left(\frac{k a_k}{4}+9b_k^2-2c_k d_k\right)
        =\frac{1}{4}\sum_{k=1}^n k a_k \;+\;9\sum_{k=1}^n b_k^2 \;-\;2\sum_{k=1}^n c_k d_k.
        \]
        Note the factor $k$ stays inside the first sum.}
        \vfill
    \end{enumerate}

\end{document}
